
\pageno=0

{\bxmplbx {\centerline {Preface}}}
\vskip .2in
\hrule height 1pt
\vskip .2in

\pageno=-1

If you're anything like me you hate to have to read large tomes devoted to 
{\it any\/} subject, even ones you like. When I see a large volume 
on a mathematical subject, it is even worse; I find even opening the 
book depressing. I'm often interested in the essential ideas of a 
subject and I don't want to wade through an encyclopedic reference. 
My experience from watching others -- as well as myself -- is that essential
ideas get lost 
in the endless details of such a reference. 

Ironically, you can see this 
in the way people {\it forget\/} subject material they've learned. 
This is not so 
strange if you think about it. If you've ever had the experience of leaving 
a town after living there a while I think you'll see what I mean.
You'll forget the names of streets after a bit. But, most likely, you will 
forget the more insignificant streets and places first. If you were to go back to this town, you 
could probably pick up the information you lost fairly easily. You 
could do this even if the town had grown or changed a fair amount. 
If instead, you were to memorize a list of street names with disconnected
snapshots of locations I think you would have a different 
experience remembering/forgetting them. You are just as likely to 
forget a major road 
or landmark as a minor one. Of course, if certain roads were set in bold in
the list
and you were told that they were more important than the others, this would 
help you remember them. What helps is to arrange the 
information in a hierarchy. In this way we prioritize the more important 
bits of information.
The idea of hierarchy occurs in  other academic areas as well. If you read a 
novel there will be major and minor themes in the text. These themes help you remember and connect the novel to 
other works. For instance, the notion of death and rebirth, or free will versus fate. 
You may forget much of the details of the novel after you have read it but 
the themes that the book is dealing with help you remember the plot. While you will forget subject 
material in this book, I would hope that
you will remember the larger themes. If exposed to calculus 
after reading this text, you will (hopefully) 
be able to recall 
the essential nature of the subject.

You may be saying ``Well of course the subject should be organized in a
hierarchy -- every text book is arranged that way.'' This is certainly true,
but in academic subjects, and mathematics in particular, there is a tradition of
creating a set of supporting concepts which build to a major result. 
This process (or series of such) leads to a crescendo with a powerful/amazing result
proclaimed. My observation watching students is that they are disconnected from
the material. The feeling one gets with mathematical instruction is that 
education is done {\it to you\/}, rather than one being an active participant in 
the exploration of the ideas one generates when attempting to solve a problem.
The feeling is one of being an audience member of a magic act. The magician does his thing and you are left to wonder how it was done.

Consequently, the material becomes a linear sequence of 
small sections that, step by step, obscure the bigger picture of the subject. 
I first saw this process
when I was a young boy. My mother would try to teach me how to balance a check 
book. It was a step by step process yet I found I couldn't remember the steps
if I hadn't done them recently. The problem was that I didn't see the 
larger purpose, it was all a jumble of steps that didn't make sense to me. 
Later on, I understood the purpose and 
could remember what to do. I was also more {\it adaptable\/} to change.
If the procedure needed to be tweaked, I could make small changes; I wasn't 
locked into doing the steps exactly the way I was told. My observation is that 
although books are organized by the hierarchy of chapters, sections, and so 
forth as any book is, they often fail to impart a useful 
{\it conceptual hierarchy\/} to the student. This is because text 
books are trying to squeeze a great deal 
of detailed information between pages. Students typically only see a blur 
of sections. Regardless of the book's hierarchy,
students' perception is often of a linear collection of sections; each as
important as the last or the next.

Today, mathematics texts are covering 
more and more 
ground at shallower and shallower depths. The argument that is often used is 
``They need to know about such and such before they take course such 
and such.''%
\footnote{\kern 0.5pt \raise 0.5ex \hbox{\dag}}%
{Strangely, physics courses are often  
six months to a year ahead of the mathematics they use.} 
I much prefer someone who has an understanding of the basics and 
can use this base to solve problems than someone who has memorized a long 
list of information for a test. To do this, I think you need to understand 
how to think about your subject. This is the intent of this book, to get 
you to see the how and the why of the {\it essential\/} concepts 
of calculus. This book is not meant to be encyclopedic; there are already 
a host of such references.

Regardless of my educational preferences, there is the looming
issue of AI which promises to do much of the detailed work in many areas of
endeavor. Going forward, we can see that there is going to be a greater 
need for humans to have a better 
understanding of the big picture of a subject. With this enhanced understanding
we can ask better questions of AI and also better direct it. 
This preface was written, in part, through a collaboration with AI in August 2026.
I have been surprised at how quickly AI has
advanced in the last few years. However, I believe that there will still be a 
need for humans to understand the big picture of calculus, and to be able to 
ask good questions of AI. This book is meant to help you develop that understanding.
To provide this understanding and remain an anti-tome book, there are a limited 
number of exercises. However, this book paired with a cheap book of exercises will provide a good understanding of the subject.
One such source of exercises is "Schaum's Outline of Calculus" by Frank Ayres and Elliott Mendelson.
The final appendix of this book was written in conjunction with AI. It provides a list of other sources for exercises; and more importantly,
a way to use AI to generate exercises and solutions. This is a new way to learn and practice calculus
and, ideally, a way to use AI to help rather than replace humans.

There is another aspect of learning that I think is underappreciated.
A student needs to feel -- at least to some extent -- that 
they {\it could have\/} thought of the ideas being presented. Not that 
they would see the answer immediately, but more that they are a 
{\it player\/} in the subject. The feeling should be something like: 
``Yeah, I think I could have done that, looking at things this way, 
given enough time.'' Without this feeling, I'm afraid all human beings 
will learn how to mechanically manipulate symbols to get an ``A'' 
on a test and never engage with the understanding.
This is similar to the way humans -- and animals -- play games. If the 
weaker animal doesn't win at least a third of the time, it loses interest in 
playing. 

How do we keep the student as a player? A key 
idea is that new material should appear as a recognizable 
{\it variant\/} of a strategy the student already possesses. 
Consider first year algebra. To solve 
$a x + b = c$ one treats the equation as a balance and plays a game of 
adding, subtracting, multiplying, and dividing to isolate $x$ on one 
side. To solve a quadratic, $a x^2 + b x + c = 0$, one manipulates 
until all the $x$'s sit inside a square, then takes a square root -- and 
is back to a problem already understood. To solve a difference equation, 
one gets all the $x$'s under a difference operator on one side and applies 
the inverse operation, summation, to get back to an algebraic equation. 
To solve a differential equation, one does the same thing with a derivative 
and its inverse, the integral. The {\it strategic move\/} is always the 
same: isolate the unknown under one operation, then apply the inverse.
What changes at each level is only {\it which\/} operation one is inverting. 
A student who sees this thread is playing the same game at every level, 
with new pieces.

The reason this works -- the reason a small number of strategies can 
carry a student from arithmetic through calculus and beyond -- is that 
mathematics reuses the same ideas at increasing levels of abstraction. 
What changes is the setting: numbers become vectors, vectors 
become functions, functions become elements of abstract spaces. 
But the game is the same. We might call this principle 
{\it same game, bigger board\/}. The student who has learned to isolate 
and invert in first year algebra is playing the same game when 
summation undoes differencing, or when integration undoes 
differentiation. The board got bigger; the move didn't change. 
Recognizing this is what keeps the student a player as the 
material grows more sophisticated.

To this end, we will introduce 
concepts through motivating examples.  This means that
the concepts may not appear in the order that you would see them in other
calculus books, but it is hoped that you will see the point of the various
definitions, theorems, and computational rules that arise. This is not to say 
that other texts don't use motivating examples. It is the case here that 
the examples will be of primary importance in creating and understanding  
the hierarchy of concepts of calculus. Major themes will be addressed straight 
away. In confronting these major themes we will see the need for a number 
of supporting ideas. Rather than build up to a major theme by 
investigating related minor themes, the minor themes will be seen as ancillary 
to our main track. This is what is meant when we say that the concepts 
may appear out of order; in fact, most of these supporting ideas will be found in the appendices.

In this regard, this text may be considered a form of 
Test Driven Development. TDD is an idea from the software engineering world 
that attempts to produce software by starting with tests first and then
writing code to make the tests pass. In a similar way, we will 
start with 
problems we wish to solve and then see what mathematical machinery is needed
to solve them. We might call this Problem Driven Development (PDD).
(If you are not a programmer, feel free to skip the analogy; the idea is
simply problems first, machinery second.)

Another organizing principle is the recurring theme of discrete versus continuous.
Here we develop a discrete as well as the traditional continuous calculus. 
This comparison is not 
done in a traditional course on calculus, but I think you will see that the 
discrete calculus also solves important problems and reinforces
concepts from the continuous calculus.
The discrete version has the
advantage of not having to deal with the technical challenges that the
continuous case forces us to consider.

More broadly, this book is organized around a small number of 
{\it strategic moves\/} that recur across all of calculus and, 
indeed, across most of the mathematics one encounters through 
an undergraduate education:
\beginEnum
\enum{{\it Isolate and Invert.\/} Get the unknown under a single
operation, then apply the inverse to peel it off. This is the thread
that connects solving linear equations with solving difference
equations and differential equations.}
\enum{{\it Approximate, then take a limit.\/} Replace a hard
continuous problem with a simple discrete one and ask what happens
as the approximation is refined. Secant lines become tangent lines;
piecewise constant sums become integrals. The discrete calculus
developed in the early chapters makes this move especially visible.}
\enum{{\it Change the representation.\/} Rewrite the problem in
a form you already know how to handle. Completing the
square turns a quadratic into a problem solvable by taking a
square root. Substitution in integration turns an unfamiliar
integral into a familiar one. The move is always the same: transform the
problem so that an existing tool applies.}
\enum{{\it Exploit linearity.\/} When an operation is linear -- it
distributes through addition and constant multiplication -- you
can break a problem into pieces, solve each piece, and add the
answers. Distributivity in arithmetic is the same structural idea
as the linearity of differentiation, integration, and summation.}
\endEnum
\medskip
\noindent These four moves cover a remarkable amount of ground. A student 
who holds these threads can walk into any course in the range from 
algebra through calculus and ask: ``Which game is this? Which move 
applies here? What is the new tool for executing it?''  
They remain a player.

We start by looking at problems dealing with the differencing and the summing 
of discrete data. In the process we find 
that differencing and summing are essentially inverses of one another. Next, 
we look at an analog of differencing and summing applied to continuous 
data; what 
are called differentiation and integration and find a similar relationship
between them.

I expect the reader to have a good understanding of algebra and algebraic notation. In particular,
we take $x\, y$ to mean $x$ times $y$; if there is confusion we will explicitly write $x * y$.
Although trigonometric
functions appear in a few sections, they are not necessary for an understanding of the rest of the material.
Finally, I hope that your experience reading this work 
gives you an understanding of calculus that will last when you leave this book.

\vskip .3in
\leftline{\it R. Scott McIntire}
\vskip .15in
\leftline{\it May 2009, revised Aug 2026}

\vfil
\supereject
